% 第 18 章。数字：checks/check_ch18.py；引文：data/snapshots/smith_facts.csv（外文史料与研究文献逐字核对见
% data/fetch/verify_web_quotes.py，《清史稿》见 data/fetch/verify_quotes.py；人口普查局的表按页面图像读，见快照元数据）；
% 数据：boe_millennium_gb_1660_1780（data/fetch/boe_millennium_gb18.py）；图：figsrc/fig_ch18.py。
\chapter{亚当·斯密：分工、积累与真实票据}\label{ch:smith}

\begin{chapterguide}
\guidequestion{国民财富从哪里来？银行应当按什么原则放贷？}
\begin{guidegoals}
\item 概括《国富论》（1776）关于分工、市场范围与资本积累的论证；
\item 解释斯密对重商主义的批评；
\item 说明真实票据论的内容，以及它后来受到批评的原因；
\item 指出斯密的哪些观点进入了后来的增长理论与货币理论。
\end{guidegoals}
\guideprereq{\chref{ch:mercantilism}\enspace\nameref{ch:mercantilism}}
\end{chapterguide}

1772 年 6 月 9 日，一位在伦敦做空东印度公司股票的银行合伙人破产，逃往法国；他做空的钱是从自己的银行借的，而这家银行又欠了苏格兰的艾尔银行（Douglas, Heron \& Company）一大笔钱\cite{rockoff2009}。6 月 12 日，一个骑马的人把消息带到爱丁堡；艾尔银行遭到挤兑，6 月 22 日被迫停止兑付自己的银行券。危机期间，爱丁堡有 15 家私人银行破产\cite{rockoff2009}。6 月 27 日，大卫·休谟（David Hume，1711—1776）写信给亚当·斯密（Adam Smith，1723—1790）：“我们这里的情形非常凄惨：接连不断的破产，信用全面丧失，猜疑没完没了”；他问斯密：“这些事件对你的理论有什么影响吗？”\cite{rockoff2009}

四年以后，斯密的《国民财富的性质和原因的研究》（\emph{An Inquiry into the Nature and Causes of the Wealth of Nations}，下称《国富论》）出版。书的第二篇第二章用了好几页讲一家“为了解救国家的困境”而设立的新银行；斯密没有点它的名，可 Rockoff 说，“他描述的是哪一家银行，是清楚的”\cite{smith1776,rockoff2009}。这本书从一句话开始：“每一个国家一年的劳动，是最初供给它一年所消费的全部生活必需品与便利品的基金”\cite{smith1776}。财富不是国库里的金银，而是一年的产出；\chref{ch:mercantilism}的重商主义者把两者混在了一起。

这一章回答导引框里的问题：国民财富从哪里来？银行应当按什么原则放贷？依次讨论：斯密写书时的英国、北美、中国与孟加拉（\secref{sec:smith-world}）；分工为什么能提高劳动生产率，什么限制了分工（\secref{sec:smith-division}）；资本怎样靠节俭积累起来，什么是生产性劳动（\secref{sec:smith-capital}）；斯密怎样批评重商主义（\secref{sec:smith-mercantile}）；纸币代替金银给国家带来了什么（\secref{sec:smith-paper}）；银行应当只贴现“真实的”票据（\secref{sec:smith-realbills}）；艾尔银行的兴起与倒闭（\secref{sec:smith-ayr}）；斯密主张怎样管银行（\secref{sec:smith-regulation}）；进步、停滞与衰退的三种国家（\secref{sec:smith-states}）。\secref{sec:smith-test} 用人口、国债、产出与银行的数据检验斯密的几个判断。《国富论》的引文据古腾堡计划的电子书，它的文本出自第三版或更晚的某一版；主要的段落在 1776 年的初版里核对过，后来增补的地方随文注明。本章引用的外文史料，中文都是本书的翻译。

\section{当时的世界}\label{sec:smith-world}

这一节交代：斯密写书的时候，英国有多富、欠了多少债；苏格兰的银行与纸币是什么样子；北美与东方又是什么样子。

\subsubsection{一个慢慢变富的国家}
按英格兰银行整理的“千年宏观数据”，大不列颠的人口 1700 年约 621 万，1776 年约 838 万；按 2013 年的价格计，人均实际 GDP 从 1,610 英镑升到 1,977 英镑，76 年里增加了 22.8\%，平均每年 0.27\%\cite{boe2017,broadberry2015}。增长这样慢，当时的人很难觉察；斯密说，那时不断有小册子出版，“声称要证明国家的财富正在迅速减少”\cite{smith1776}。

国家欠的债却增长得很快。斯密列举 1688 年以来的“四次耗资巨大的对法战争”，分别开始于 1688、1701、1742 与 1756 年，另有 1715 年与 1745 年的两次叛乱；他说，四次对法战争里，国家“欠下了超过 1.45 亿英镑的债”，连同战争引起的其他非常开支，“总数不会少于 2 亿英镑”\cite{smith1776}。按千年数据里的国债（面值），1700 年是 1,410 万英镑；1763 年七年战争结束\cite{rockoff2009}，国债是 1.342 亿英镑，算上定期年金的资本价值是 1.442 亿英镑，与斯密的数字相近；1776 年是 1.366 亿英镑，算上年金约为 GDP 的 97.5\%\cite{boe2017,mitchell1988}。斯密的结论是，政府的浪费“无疑延缓了英格兰走向财富与改良的自然进程，却没能让它停下来”\cite{smith1776}。图~\ref{fig:smith_debt} 画出了这两条曲线。

\subsubsection{苏格兰：纸做的钱}
苏格兰有两家“公共银行”：苏格兰银行 1695 年由议会法案设立，皇家银行 1727 年由国王特许设立\cite{smith1776}。斯密估计，1707 年合并以前，苏格兰流通的金银“不会少于一百万英镑”，其中送进苏格兰银行重铸的银币是 411,117 英镑 10 先令 9 便士；他写书的时候，苏格兰的全部通货“不会少于两百万英镑”，其中金银“很可能不到五十万英镑”，其余都是银行券\cite{smith1776}。苏格兰的银行券小到 20 先令一张；银行还发明了“现金账户”，凭两个有信用、有地产的保证人，给一个人开出一定额度的信用，“比如两三千英镑”\cite{smith1776}。斯密听人说，格拉斯哥的银行设立以后，那里的贸易大约十五年就翻了一番，苏格兰的贸易自两家公共银行设立以来增加到四倍以上；贸易是不是真在这么短的时间里增长了这么多，“我不敢说知道”，即使真是这样，单靠银行也解释不了这么大的效果；不过他认为，苏格兰的贸易与产业这期间增加了很多，银行出了不少力，这一点无可怀疑\cite{smith1776}。

\subsubsection{北美与东方}
斯密拿来与英国对比的有三个地方。北美的英国殖民地增长最快：在大不列颠与欧洲多数国家，居民“据认为不到五百年不会翻一番”，北美的英国殖民地“二十或二十五年就翻一番”\cite{smith1776}。按美国人口普查局根据殖民地官员的报告整理的估计，北美殖民地的人口 1700 年约为 25.1 万，1770 年约为 214.8 万\cite{census1975}。中国“似乎长期停滞”：五百多年以前马可·波罗（Marco Polo，1254—1324）对它的耕作、工业与人口的描述，与今天的旅行者“几乎一样”\cite{smith1776}。孟加拉在衰退：斯密描述了一个衰退的国家，工资降到最低，饥荒与死亡从最穷的人蔓延开来，然后说，“这大概就是孟加拉……目前的状况”；在一个土地肥沃、谋生本不太难的国家，如果“一年里饿死三四十万人”，就可以断定维持劳动贫民的基金正在迅速衰减\cite{smith1776}。这三个地方是\secref{sec:smith-states}的题目。

\section{分工与市场的范围}\label{sec:smith-division}

这一节回答：国民财富从哪里来？分工为什么能提高劳动生产率？什么限制了分工？

\subsubsection{两个因素}
斯密说，一个国家拿来消费的产出，与消费它的人数之比，“取决于两种不同的情况：第一，劳动通常运用得是否熟练、灵巧、有判断；第二，从事有用劳动的人数与不从事有用劳动的人数之比”\cite{smith1776}。设一年的产出为 $Y$，人口为 $N$，从事有用劳动的人数为 $L$，这句话就是一个恒等式：
\begin{keyeq}
\frac{Y}{N}=\frac{Y}{L}\cdot\frac{L}{N} .\label{eq:smith-twofactors}
\end{keyeq}
人均产出等于劳动生产率 $Y/L$ 乘以从业的比例 $L/N$。斯密说的是从业者与不从业者之比 $L/(N-L)$，它与 $L/N$ 同升同降，用 $L/N$ 写出来更简单。恒等式本身不说明任何因果，它的用处是把问题分开：人均产出的增长，要么来自生产率，要么来自更多的人去做事。斯密说，供给的多寡“似乎更多地取决于前一个情况，而不是后一个”\cite{smith1776}。他的证据是一个对比：在以渔猎为生的民族里，能劳动的人几乎都在劳动，却穷得有时不得不遗弃婴儿、老人与久病的人；在文明而繁荣的国家里，许多人根本不劳动，全社会的产出却多到人人常常得到充分的供给\cite{smith1776}。用式~\eqref{eq:smith-twofactors} 来说，$L/N$ 不可能超过 1，$Y/L$ 却没有这样的上限。

\subsubsection{别针工场}
斯密用一个别针工场说明分工。一个没有学过这一行的工人，“尽最大的努力，大概一天也做不出一枚别针，肯定做不出二十枚”；分了工，做一枚别针要经过“大约十八道不同的工序”\cite{smith1776}。斯密说：“我见过一家这样的小工场，只雇了十个人”，他们一天“能做出大约十二磅别针”；一磅中等大小的别针有四千多枚，所以这十个人“一天能做出四万八千多枚别针”，每人折合四千八百枚。如果各干各的，他们“肯定每人一天做不出二十枚，也许连一枚也做不出；就是说，肯定不到他们现在能做的二百四十分之一，也许不到四千八百分之一”\cite{smith1776}。

斯密把这样大的差别归于三种情况：每个工人熟练程度的提高；节省了“从一种工作转到另一种工作时通常损失的时间”；以及发明了许多机器，它们“便利和简化了劳动，使一个人能做许多人的工作”\cite{smith1776}。从增长的角度看，第三条最要紧，因为它可以一直积累下去：分工让一个人整天只做一件事，他就容易想到用机器去做。

\begin{example}{别针与两个因素}\label{ex:smith-pins}
（1）用斯密的数字验算：十个人一天做 12 磅别针，每磅 4,000 枚，每人每天做多少枚？与一个人每天至多 20 枚、至少不到 1 枚相比，是多少倍？（2）教学假设：一个国家从事有用劳动的人占 40\%。按式~\eqref{eq:smith-twofactors}，只靠提高这个比例，人均产出最多能增加到原来的几倍？如果比例从 40\% 升到 50\%，人均产出增加多少？（3）大不列颠 1700—1776 年的人均产出每年增长 0.27\%。按这个速度，人均产出翻一番要多少年？
\solution （1）$12\times4{,}000=48{,}000$ 枚，每人 4,800 枚；$4{,}800/20=240$，即斯密的“二百四十分之一”，与 1 枚相比是 4,800 倍。（2）$L/N$ 最多是 1，人均产出最多增加到 $1/0.4=2.5$ 倍；从 40\% 升到 50\%，人均产出增加 $0.5/0.4-1=25\%$。生产率没有这样的上限，所以斯密更看重它。（3）$\ln2/\ln1.0027\approx257$ 年。按这个速度，一个人一辈子几乎看不出变化，难怪有人写小册子说国家在变穷。
\end{example}

\subsubsection{市场的范围}
分工有一个限度。斯密说，分工起于交换的能力，“所以分工的程度，总要受这种能力的大小限制，换句话说，受市场范围的限制”\cite{smith1776}。他举了一个钉匠的例子：一个钉匠“每天做一千枚钉子，一年做三百个工作日，一年就做三十万枚钉子”；可是在苏格兰高地的偏远地方，“一年也不可能卖出一千枚，也就是一天的活儿”\cite{smith1776}。那里的人只能什么都自己做，或者一个铁匠兼做许多种活儿。

市场的范围又取决于运输。斯密比较了陆运与水运：一辆宽轮的大车，“由两个人照料，八匹马拉着，大约六个星期，在伦敦与爱丁堡之间运去又运回将近四吨货物”；而一条“由六个或八个人驾驶”的船，在伦敦与利斯两港之间航行，常常运去又运回两百吨货物，相当于“五十辆宽轮大车，由一百个人照料，四百匹马拉着”\cite{smith1776}。所以最早的文明都在水边。斯密问：“伦敦与加尔各答之间，什么货物负担得起陆运的费用？”他还说，“在孟加拉，恒河与其他几条大河形成了许多可以通航的水道，就像尼罗河在埃及一样”；中国的东部各省，几条大河的支流也形成了“许许多多的运河”\cite{smith1776}。图~\ref{fig:smith_division} 把别针工场与运输的数字画在一起。

\fig{smith_division}{分工与市场的范围。(a) 看三个点相差多少个数量级：一个没有学过这一行的工人一天做的别针，斯密说“也许连一枚也做不出”“肯定做不出二十枚”；分了工的工场里每人每天 4,800 枚（十个人一天 12 磅，每磅四千多枚）。横轴为对数刻度。(b) 看陆运与水运运送同样的 200 吨货物要多少人与马：五十辆宽轮大车要 100 个人与 400 匹马，一条船只要 6 到 8 个人。(c) 一个钉匠一年能做三十万枚钉子，苏格兰高地的偏远地方一年卖不出一千枚：市场太小，分工就停在那里。数据：斯密《国富论》第一篇第一章、第三章。}{fig:smith_division}

\begin{example}{一个钉匠需要多大的市场}\label{ex:smith-market}
（1）运送 200 吨货物，陆运要多少辆车、多少人、多少匹马？用的人是水运的几倍？（2）教学假设：一户人家一年用 50 枚钉子。一个专职的钉匠要有多少户人家买他的钉子，才能整年开工？一个只有 15 户人家的高地村子，一年用多少枚？够他做几天？
\solution （1）每辆车运 4 吨，$200/4=50$ 辆，$50\times2=100$ 个人，$50\times8=400$ 匹马，与斯密的数字一致；人数是水运的 $100/8=12.5$ 到 $100/6\approx16.7$ 倍，还没算马。（2）$300{,}000/50=6{,}000$ 户；15 户一年用 $15\times50=750$ 枚，不到一天的产量（1,000 枚）。斯密说高地的钉匠一年卖不出一千枚，正是这个意思：在那里，钉匠不能只做钉子。
\end{example}

\section{资本：生产性劳动与节俭}\label{sec:smith-capital}

这一节回答：什么样的劳动是“生产性的”？资本怎样积累？国家的浪费会不会让国家变穷？

\subsubsection{生产性劳动}
第二篇第三章开头说：“有一种劳动，加在对象上能增加它的价值；另一种没有这种作用。前者因为生产价值，可以叫作生产性劳动，后者叫作非生产性劳动。”\cite{smith1776}他的例子是：“一个人雇用许多工匠就变富，养活许多仆人就变穷。”不过，“后者的劳动也有它的价值，也和前者一样应得到报酬”；区别在于，工匠的劳动“固定并实现在某种特定的对象或可以出售的商品上，这种商品在劳动过去以后至少还能保存一段时间”\cite{smith1776}。按这个标准，“君主，以及在他手下任职的全部司法与军事官员，整个陆军与海军，都是非生产性劳动者”；他们的服务“无论多么光荣、多么有用、多么必要”，都不生产以后能换来同量服务的东西。同属一类的还有“牧师、律师、医生、各种文人；演员、小丑、音乐家、歌剧演员、歌剧舞蹈演员等等”\cite{smith1776}。这个标准比重农学派的宽：魁奈只把农业劳动算作生产性的，把工匠、制造业者与商人归入不生产阶级（\secref{sec:merc-tableau}），斯密认为这正是那个体系的根本错误。

\begin{compare}{两种“生产”}{斯密的生产性劳动}{今天的 GDP}
\cmprow{什么算产出}{固定在可以出售、能保存一段时间的物品上的劳动}{一年里生产的全部最终货物与服务（\chref{ch:accounts}）}
\cmprow{仆人、律师、音乐家}{非生产性：服务随做随消失}{他们的服务计入 GDP}
\cmprow{君主、陆军与海军}{非生产性劳动者}{政府提供的服务也计入 GDP}
\cmprow{区分的用意}{哪些开支能变成资本、雇用更多生产性劳动}{一年的产出与收入有多少}
\cmprow{斯密怎样评价非生产性劳动}{“也有它的价值，也应得到报酬”}{——}
\cmpsum{斯密的区分不是一条统计口径，而是一个关于积累的论断：收入花在雇用工匠上，能留下资本，花在雇用仆人上，就留不下来。今天的国民账户把两种服务都算作产出，积累的问题交给储蓄与投资（式~\eqref{eq:accounts-balances}）。}
\end{compare}

\subsubsection{节俭}
斯密说：“资本因节俭而增加，因奢侈与行为失当而减少”；“勤劳不是资本增加的直接原因，节俭才是”\cite{smith1776}。省下来的钱并没有退出流通：“每年省下来的，与每年花掉的一样，有规律地被消费掉，而且几乎在同一时间；只是消费它的是另一群人”，即用这笔钱雇来的生产性劳动者。“消费是一样的，消费者不同。”\cite{smith1776}所以，“每个浪费的人都像是公众的敌人，每个节俭的人都是公众的恩人”\cite{smith1776}。

人为什么会节俭？斯密说，促使人储蓄的，“是改善自身境况的愿望；这种愿望虽然一般是平静的、不动感情的，却从娘胎里就跟着我们，直到进坟墓都不离开”\cite{smith1776}。他相信多数人在多数时候是审慎的：经商的人里破产的，“也许不比千分之一多多少”\cite{smith1776}。斯密由此回到式~\eqref{eq:smith-twofactors}：一国一年的产出，“除了增加生产性劳动者的人数，或者提高原有劳动者的生产力，没有别的办法增加它的价值”\cite{smith1776}；前者要更多的资本来雇人，后者要更多的资本来买机器、改进分工。

\begin{pitfall}[省下来的钱一定会被花掉]
斯密说储蓄“同样被消费掉”，有时被读成一条定律：储蓄不会减少总需求，所以不会有普遍的生产过剩。斯密的论证有一个前提：省下的收入马上被当作资本，用来雇用生产性劳动者。如果省下的钱被存着不用，这个前提就不成立；斯密自己也把闲置的现金叫作“死的资财”（\secref{sec:smith-paper}）。他为这个前提给的理由在第二篇第一章：在大体安全的国家里，有常识的人都会把能支配的资财用来享用或者谋利，只有在怕被掠夺的地方，人们才把资财埋起来\cite{smith1776}。这是一个关于个人动机的论证；许多人同时多储蓄时，储蓄是否总会变成投资，是\chref{ch:say}讨论的“普遍过剩”之争与\chref{ch:generaltheory}的节俭悖论的问题。
\end{pitfall}

\subsubsection{公家的浪费}
斯密对私人很乐观：“大国从来不会因私人的浪费与失当而变穷，虽然有时会因公家的浪费与失当而变穷。”\cite{smith1776}他对君主与大臣毫不客气：他们自己“永远、毫无例外，是社会上最大的浪费者”\cite{smith1776}。\secref{sec:smith-world}引了他对四次对法战争的估计：超过 1.45 亿英镑的债，总共不少于 2 亿英镑的开支。这些钱如果没有被借去打仗，大部分会被用来雇用生产性劳动者；可是，斯密说，政府的浪费只是延缓了进步，“却没能让它停下来”。按千年数据，1700—1763 年，国债（面值）从 1,410 万英镑增加到 1.342 亿英镑，人均实际 GDP 也从 1,610 英镑升到 1,945 英镑，增加了 20.8\%\cite{boe2017,broadberry2015}。

\fig{smith_debt}{英国的国债与人均产出，1660—1780 年。(a) 看曲线在四条竖线以后怎样上一个台阶：实线是国债的面值（百万英镑），褐色虚线另加定期年金的估计资本价值；竖线是斯密说的四次对法战争开始的年份（1688、1701、1742、1756），水平虚线是他说的“超过 1.45 亿英镑”。(b) 看两条线在一百年里升了多少：大不列颠（1700 年起）与英格兰的人均实际 GDP，按 2013 年价格计（英镑）。数据：英格兰银行《千年宏观数据》表 A21、A29（Broadberry 等 2015；Mitchell 1988）。}{fig:smith_debt}

\section{财富不是金银}\label{sec:smith-mercantile}

这一节回答：斯密怎样批评重商主义？他的批评留下了什么问题？

\secref{sec:merc-mun}讲过孟怎样用贸易差额解释国家的财富。斯密在第四篇第一章里说：“要认真去证明财富不在于货币或金银，而在于货币能买到的东西，金银的价值只在于能买东西，这未免太可笑了。”\cite{smith1776}货币诚然是国民资本的一部分，可它“通常只是一小部分，而且总是最不能生利的部分”\cite{smith1776}。人们常抱怨钱少，斯密说：“引起普遍抱怨钱少的，不是金银的缺乏，而是这些人借钱的困难，以及他们的债主收债的困难”\cite{smith1776}。他把因果倒了过来：一个国家的产出在衰退，它的金银就会流出，“在这种情况下，金银的出口不是衰落的原因，而是衰落的结果”\cite{smith1776}。

重商主义者常把国家比作商人：商人要卖出比买进多，才能积攒钱财。斯密不接受这个比喻：“一个商人仓库里货物很多，有时会因为不能及时卖掉而破产，一个国家却不会遭遇同样的意外”\cite{smith1776}。这句话后来成了一场大争论的起点：整个国家的货物会不会普遍卖不出去？那是\chref{ch:say}的题目。

斯密反对用关税与禁令保护国内产业，最常被引用的一句话就出现在这里。一个人宁愿把资本用在国内而不是国外，“他想的只是自己的安全”；他把资本用在产值最大的地方，“他想的只是自己的得益；在这里，像在许多别的场合一样，他被一只看不见的手引导着，去促进一个并非他本意的目的”\cite{smith1776}。接着，斯密说：“他追求自己的利益，往往比真心想促进社会利益时更有效地促进了社会利益。我从没听说那些假装为公益而做生意的人做过多少好事。”\cite{smith1776}这一段在初版的第 2 卷里就有\cite{smith1776}。

\begin{compare}{财富是什么}{孟（\chref{ch:mercantilism}）}{斯密}
\cmprow{财富是什么}{金银：国家的财宝}{一年的劳动产出}
\cmprow{怎样增加}{贸易顺差，金银流入}{分工与资本积累：更多的生产性劳动者，更高的生产力}
\cmprow{货币的作用}{财富本身}{流通的轮子；闲置的金银是“死的资财”}
\cmprow{国家与商人}{国家像商人，要卖得比买得多}{国家不像商人，不会因为货物卖不出去而破产}
\cmprow{政策}{管制贸易，追求顺差}{取消管制，让资本流向产值最大的地方}
\cmpsum{孟问的是一个国家怎样积攒金银，斯密问的是一个国家一年能生产多少。斯密的问法留下了一个新问题：一国的产出会不会因为需求不足而卖不出去。}
\end{compare}

\section{纸币：空中的车道}\label{sec:smith-paper}

这一节回答：用纸币代替金银，国家得到什么？一家银行多发了纸币，会怎样？斯密在同一章还讨论了北美殖民地政府发行的纸币，\secref{sec:money-functions}与\exref{ex:money-smith}已经讲过，这里不再重复。

\subsubsection{一个新的轮子}
斯密把金银比作流通的轮子：它花费很大，自己却不生产任何东西。商人为应付随时之需留在身边的现款，“是死的资财，只要还留在这种状态，对他、对他的国家都不生产任何东西”；一国流通的金银也是这样\cite{smith1776}。用纸币代替金银，“流通就由一个新的轮子来推动，造它、养它都比旧轮子便宜”\cite{smith1776}。他的算术是这样的：一个银行家发出十万英镑的银行券，通常只要留两万英镑的金银，“就足以应付随时的兑现要求”；于是，“八万英镑的金银可以这样从国家的流通中节省下来”，流通只用原来所需金银的“五分之一”就能进行\cite{smith1776}。

斯密再往下推：设一个国家流通的货币原来是一百万英镑金银，银行发出一百万英镑的银行券，“在各自的金库里留出二十万英镑应付随时的兑现要求”，流通中就有八十万英镑的金银加一百万英镑的银行券，共一百八十万英镑\cite{smith1776}。可是，一国的产出没有变，“流通的渠道，如果允许我这样说，仍然和以前一样”；于是，“八十万英镑必然溢出”。溢出的不会是纸币：“纸币不能出国，因为离开发行它的银行、离开能依法要求兑付的国家，它在日常支付中不会被接受”\cite{smith1776}。溢出的是金银，它们被送到国外，换回货物。斯密认为，这些货物大部分是原料、工具与食物，用来雇用更多的劳动；所以，金银被纸币替换下来以后，即使只有“其余五分之四里的大部分”被加进维持劳动的基金，“也必然大大增加这种劳动的数量，从而大大增加一年的土地与劳动产出的价值”\cite{smith1776}。

这个“大大增加”有多大？设流通的金银为 $M$，银行留作准备的比例为 $\rho$，换回来的资本的利润率为 $r$，产出 $Y$ 增加的比例是
\begin{keyeq}
\frac{\Delta Y}{Y}=\frac{M}{Y}\,(1-\rho)\,r .\label{eq:smith-saving}
\end{keyeq}
三个因子分别是货币占产出的比例、被替换下来的金银的份额，以及这笔金银换成资本以后的回报。斯密自己说，货币占一年产出的比例很难确定，“不同的作者估计为五分之一、十分之一、二十分之一与三十分之一”\cite{smith1776}。Rockoff 按斯密的例子做了计算：设贩运贸易的利润率是每年 10\%，八十万英镑每年能赚八万英镑；按货币占产出五分之一计，银行带来的增加是产出的 1.60\%，按最不利的三十分之一计，只有 0.27\%；“即使是较大的数字 1.60\%，也是比较小的”\cite{rockoff2009}。

Rockoff 猜想，斯密心里想的是一种回报高得多的投资：种子。设要用一万吨种子才能收获十万吨粮食；如果用金银从国外买来粮食，消费不变，多留出一万吨种子，收成就会翻一番。Rockoff 说：“在这里，我认为，斯密打了个盹”（英语里借荷马也会打盹的说法，指一时疏忽）；斯密“暗中假定农业投资的回报率非常高”，在这个例子里是 900\%\cite{rockoff2009}。

\begin{example}{纸币节省了多少}\label{ex:smith-saving}
（1）按式~\eqref{eq:smith-saving}，$\rho=1/5$、$r=10\%$，货币占产出的比例取斯密列举的四个数，产出各增加多少？（2）要让产出增加 10\%，按五分之一与三十分之一，利润率 $r$ 各要多高？（3）Rockoff 的例子里，一万吨种子多收九万吨粮食（十万吨减去种子本身），回报率是多少？
\solution （1）$\tfrac15\times\tfrac45\times10\%=1.6\%$；十分之一为 0.8\%，二十分之一为 0.4\%，三十分之一为 $0.27\%$。前后两个数正是 Rockoff 的 1.60\% 与 0.27\%。（2）$10\%/(0.2\times0.8)=62.5\%$；$10\%/(\tfrac1{30}\times0.8)=375\%$。斯密说的“大大增加”，要求这样高的回报。（3）$(100{,}000-10{,}000)/10{,}000=900\%$，正是 Rockoff 的数字。他的意思是：只有假定这样高得离谱的回报，斯密的“大大增加”才说得通。纸币的好处是真实的，但按正常的利润率算，它一次性地给产出加上不到两个百分点，不是增长的源泉。
\end{example}

\begin{intuition}[空中的车道与代达罗斯的翅膀]
斯密说，一国流通的金银，好比一条大路，“它把全国的草与谷物运到市场上，自己却连一根草、一粒谷都不生产”；银行用纸币代替金银，就像给国家开辟了“一种空中的车道”，让国家把许多大路改成好的牧场与麦田\cite{smith1776}。他马上又警告说，国家的商业与工业“这样悬在纸币的代达罗斯翅膀上”，不能像“在金银的坚实地面上行走”时那样安全；比如“一场失败的战争”，敌人占领了首都，也就占领了支撑纸币信用的那笔财宝，整个流通都靠纸币的国家就会陷入大得多的混乱\cite{smith1776}。Rockoff 提醒读者，代达罗斯“用蜡把羽毛粘成翅膀”，他自己用得很谨慎，他的儿子伊卡洛斯却飞得离太阳太近，蜡化了，坠入海中\cite{rockoff2009}。Rockoff 说，在一份据说写于 1763 年 4 月以前的草稿里，已经有“空中的车道”，却没有代达罗斯的翅膀；1776 年的初版两者都有\cite{rockoff2009,smith1776}。他说，这个比喻用来形容 1772 年的危机，“或者说 2008 年的危机”，都再好不过\cite{rockoff2009}。
\end{intuition}

\subsubsection{多发的纸币会回来}
一家银行能不能随意多发纸币？斯密说不能，只要纸币随时可以兑换金银。设某地的流通“正好能吸收四万英镑”的纸币，银行为应付兑现，“必须随时在金库里留一万英镑的金银”。如果它想发出四万四千英镑，“超过流通容易吸收与使用的那四千英镑，几乎一发出去就会回到它手里”\cite{smith1776}。设流通能吸收的纸币为 $\bar N$，银行发出 $N>\bar N$（这里的 $N$ 指银行券，与式~\eqref{eq:smith-twofactors} 的人口无关），准备的比例仍为 $\rho$，它要留的金银是
\begin{keyeq}
R=\rho\,\bar N+(N-\bar N) .\label{eq:smith-reflux}
\end{keyeq}
多发的部分会全部回来兑现，所以要按一比一准备。斯密说，这家银行“要留的不只是一万一千英镑，而是一万四千英镑”；它从多发的四千英镑上“得不到任何利息”，却要承担不断收集四千英镑金银的全部费用\cite{smith1776}。

收集金银是要花钱的。斯密说，为了应付兑现，英格兰银行多年来每年要铸造八十万到一百万英镑的金币，“平均约八十五万英镑”；它以每盎司 4 英镑的高价买进金块，铸成金币时只按每盎司 3 英镑 17 先令 10 便士半发出，“这样损失二点五到三个百分点”\cite{smith1776}。苏格兰的银行托伦敦的代理人收集现金，费用“很少低于一点五或两个百分点”，用大车运到苏格兰，承运人保险又要“四分之三个百分点，即每一百英镑十五先令”\cite{smith1776}。斯密的结论是：“英格兰与苏格兰两地的一些大胆的投机者过度经营，是纸币过度流通的根本原因”\cite{smith1776}。

\begin{example}{多发四千英镑}\label{ex:smith-excess}
（1）用式~\eqref{eq:smith-reflux} 验算斯密的数字：$\bar N=40{,}000$，准备一万英镑，$N=44{,}000$。按同样的准备比例，发 44,000 英镑要留多少？实际要留多少？（2）每盎司按 4 英镑买进，铸成 3 英镑 17 先令 10 便士半发出（1 英镑 = 20 先令，1 先令 = 12 便士），每发出一英镑金币损失百分之几？英格兰银行每年平均铸 85 万英镑，一年损失多少？
\solution （1）$\rho=10{,}000/40{,}000=1/4$，按同样的比例要留 $44{,}000/4=11{,}000$ 英镑，即斯密的“一万一千英镑”；实际要留 $R=10{,}000+4{,}000=14{,}000$ 英镑。（2）$3+17/20+10.5/240=3.89375$ 英镑；$4/3.89375-1\approx2.73\%$，在斯密说的二点五到三个百分点之间；一年约 $850{,}000\times2.73\%\approx2.3$ 万英镑。
\end{example}

\subsubsection{纸币会不会抬高物价}
斯密转述了一种看法，没有说是谁的：“有人说，纸币的增加增加了整个通货的数量，因而降低了它的价值，必然抬高商品的货币价格。”他的回答是：“从流通中取走的金银，总是等于加进来的纸币，所以纸币不一定增加整个通货的数量。”\cite{smith1776}他举苏格兰为例：1759 年流通着十先令与五先令的小额银行券，纸币比现在多，“苏格兰的粮食却从没有那样便宜过”；1751 年与 1752 年，“休谟先生出版他的《政治论丛》的时候”，紧接在苏格兰纸币大量增加之后，粮价确实上涨了，可那“大概是由于年成不好，而不是由于纸币的增加”\cite{smith1776}。斯密的论断有一个条件：纸币随时可以兑现。他说，如果纸币的立即兑付“在任何方面取决于发行人的善意”，“情形就不一样了”\cite{smith1776}。

二十多年以后，亨利·桑顿（Henry Thornton，1760—1815）指出了这个论证的一个漏洞。斯密认为，有了纸币以后能流通的纸币，与没有纸币时会流通的金币，“数量相同”\cite{thornton1802}；可是，同一种钱流通得有快有慢，“同一张钞票可以一天完成十笔支付，也可以十天完成一笔支付”\cite{thornton1802}。流通的“渠道”不变，钱的数量仍然可以变。桑顿的论证与他对真实票据论的批评，是\chref{ch:bullion}的题目。

\section{银行应当贴现什么}\label{sec:smith-realbills}

这一节回答：银行应当按什么原则放贷？“真实的”票据与“循环的”票据有什么不同？

\subsubsection{只借他本来要闲置的钱}
斯密的原则是：“银行可以适当地贷给一个商人或任何企业主的，不是他做生意的全部资本，甚至不是资本中相当大的一部分，而只是他本来不得不闲置在身边、留作现款以应付随时之需的那一部分。”\cite{smith1776}一张\textbf{真实票据}正好满足这个条件：“一家银行给一个商人贴现一张真实的汇票，它由一个真实的债权人开出，以一个真实的债务人为付款人，到期时真的由这个债务人付清；银行只是预付给他一部分价值，这部分价值他本来不得不闲置在身边、留作现款以应付随时之需。”\cite{smith1776}票据到期，债务人付款，钱回到银行。斯密说，这样一家银行的金库就“像一个水池，虽然不断有水流出，也不断有同样多的水流进”\cite{smith1776}。

现金账户也可以照这个原则管理：银行看一个客户在一段短时期内，“比如四个、五个、六个或八个月”，还回来的钱是否等于借出去的钱\cite{smith1776}。斯密说，这样管理，除了几乎省下补充金库的全部额外费用，还有两个好处。其一，银行看得见借款人的生意：“一家把钱借给大概五百个不同的人、注意力又不断被各种事情占据的银行，除了自己的账簿提供的情况以外，对大部分债务人的行为与境况，不可能有经常的了解。”\cite{smith1776}还款按时、足额，账簿就说明了一切。其二，银行不会发出多于流通容易吸收的纸币：一个客户的还款大体等于借款，借给他的纸币就没有超过他本来要留在身边应付随时之需的金银\cite{smith1776}。

银行不应当贷给企业主的是\textbf{固定资本}：采矿的人用来“开凿矿井、安装抽水机械、修筑道路与车道”的资本，或者改良土地的资本，回收得慢；这样的钱“不应当向银行借，而应当凭债券或抵押，向那些打算靠利息生活的私人借”\cite{smith1776}。

\subsubsection{循环开票}
斯密说，二十五年多以来，苏格兰的银行发出的纸币一直“完全等于、甚至略多于国家的流通容易吸收与使用的数量”\cite{smith1776}。借款人还想借得更多：他们以为银行只要花“几令纸”的成本，就能把信用扩大到他们想要的任何数额，于是抱怨银行董事“目光短浅、胆小怯懦”\cite{smith1776}。银行不肯再扩大信用，借款人找到了另一个办法：开出没有货物交易做后盾的汇票，到期时再开一张新的去还旧的。这是“濒临破产的不幸商人有时采用的办法”，在“上次战争期间，贸易的高利润诱人过度经营”时，据说做得非常多\cite{smith1776}。斯密说：“这房子是危楼，一个疲倦的旅客对自己说，它不会撑得太久；可它今晚倒下来的可能不大，我就冒险在里面睡一晚。”\cite{smith1776}

他的例子是：“爱丁堡的商人 A 开一张汇票给伦敦的 B，两个月后付款。”\cite{smith1776}A 在爱丁堡把汇票贴现，拿到钱；汇票快到期时，B 再开一张两个月的汇票给 A，加上利息与佣金，拿去贴现，用来付第一张；如此往复，有时持续好几年。斯密算了它的成本：“利息是每年 5\%，每开一次票，佣金从不低于 0.5\%。佣金一年要付六次以上，A 用这个办法筹到的钱，一年的成本必然超过 8\%”；“这种做法叫作靠流通筹钱”\cite{smith1776}。设每年开票 $n$ 次，每次的佣金为 $k$，利率为 $i$，一年的成本是
\begin{keyeq}
c=i+n\,k .\label{eq:smith-redraw}
\end{keyeq}
斯密说，在一个大部分商业项目的普通利润“据认为在 6\% 到 10\% 之间”的国家，成本超过 8\% 的钱，只有极其幸运的项目付得起\cite{smith1776}。

更要紧的是，这种票据不会流回银行。斯密说，用新票付旧票，“这种付款完全是虚构的”；靠这种循环的汇票从银行金库流出去的水，“从来没有被真正流进来的水补上”\cite{smith1776}。两个人总在同一家银行贴现，银行家马上就能看出他们在做什么；可是一大群投机者轮流在不同的银行贴现、互相帮忙开票，就能让人“尽可能难以分辨真实的汇票与虚构的汇票”\cite{smith1776}。银行发现以后，一天天收紧贴现，这些投机者大为恼怒；斯密说，“他们把自己的困境叫作国家的困境”\cite{smith1776}。图~\ref{fig:smith_banks} 把只贴现真实票据的银行与下一节的艾尔银行放在一起：前者放出去的银行券随票据到期流回来，后者放出去的贷款回收很慢，只能在伦敦开新票还旧票。

\begin{example}{靠流通筹钱的成本}\label{ex:smith-redraw}
（1）用式~\eqref{eq:smith-redraw} 算斯密的例子：两个月的汇票，一年开六次，佣金每次 0.5\%，利率 5\%。（2）后来的版本在这里加了一条注：战争快结束时，爱丁堡对伦敦的汇价常常贴水 3\%，有一种做法一年至少重复四次，每次另付至少 0.5\% 的佣金\cite{smith1776}。一年的成本是多少？（3）桑顿举的例子是两个月的票连开三次，他说，这种交易“实质上显然等于 A 与 B 以两人的共同担保借了这笔钱六个月”，而且常常轮流为对方做，互相付佣金\cite{thornton1802}。如果两人轮流，各开一半的票，每人一年净付的佣金是多少？
\solution （1）$c=5\%+6\times0.5\%=8\%$，斯密说“超过 8\%”，因为有时开票更频繁、佣金更高，或者要为以前的利息与佣金付复利。（2）$4\times(3\%+0.5\%)=14\%$，正是注里的“至少 14\%”；这条注初版没有（\secref{sec:smith-test}）。（3）一半的时候是为 A 筹钱，A 付佣金给 B；另一半的时候是为 B 筹钱，B 付佣金给 A。每人付出的佣金与收进的大致相等，净额为零，只剩下 5\% 的利息（贴现）。桑顿据此说，这种做法常常“比斯密想象的”便宜得多\cite{thornton1802}。斯密的论点不全靠成本：即使利息只有 5\%，这种票据背后也没有会到期付款的货物，钱不会流回银行。
\end{example}

\fig{smith_banks}{两种银行的资产负债表。(a) 看左边的票据怎样不断到期、钱怎样流回：一家只贴现真实票据的银行，资产是金银准备与两个月内到期的商业票据，负债是银行券；票据到期，债务人付款，流出的水等于流进的水（斯密的“水池”）。(b) 看贷款的箭头为什么只朝外：艾尔银行 1772 年 6 月停止兑付时，银行券约 20 万英镑，在伦敦开出的汇票超过 60 万英镑，贷出超过 80 万英镑；贷款用于改良土地与续借的循环汇票，回收很慢，银行只能不断在伦敦开新票还旧票，“像派人提着水桶到几英里外的井里打水”。示意图，数字据斯密《国富论》第二篇第二章。}{fig:smith_banks}

\begin{compare}{真实票据论与货币数量论}{真实票据论}{货币数量论}
\cmprow{回答的问题}{一家银行该贷给谁、贷多少}{物价水平由什么决定}
\cmprow{主张}{只贴现真实票据，纸币到期就会流回，不会过多}{物价随货币的数量变化（\chref{ch:silver}的式~\eqref{eq:silver-qty}）}
\cmprow{物价的锚}{斯密：纸币随时可以兑换金银}{货币的数量}
\cmprow{纸币不兑现时}{票据的名义价值随物价上涨，发行没有上限（Mints）}{仍然成立}
\cmprow{斯密的立场}{审慎经营的规则（Laidler）}{否认可兑现的纸币会抬高物价：1759 年与 1751—1752 年}
\cmpsum{两种理论回答的不是同一个问题。在金银本位下，真实票据只是一条防止银行冒险的规则；把它当作控制货币总量的规则，就要另找一个锚。}
\end{compare}

\begin{pitfall}[真实票据论是一条控制货币数量的规则]
Mints 把斯密称为“真实票据学说第一个彻底的阐述者”\cite{rockoff2009}。Mints 的批评是：在通胀的经济里，送到银行贴现的票据的名义价值会随物价上涨，“所以即使每一家银行、包括中央银行都遵循真实票据原则，通胀也可以持续下去”\cite{rockoff2009}。Rockoff 说，他在芝加哥读研究生的 20 世纪 60 年代，这一学说总被称为“谬误的真实票据学说”\cite{rockoff2009}。但是 Rockoff 也说，是后来的理论家与决策者“犯了把真实票据当作货币政策规则的错误”；按 Laidler 的看法，斯密从没有犯这个错误：“在斯密的模型里，物价水平由金银本位来锚定；真实票据是为了防止不审慎的银行经营。”\cite{rockoff2009}Rockoff 还说，斯密没有打算让真实票据成为国家建立的管制结构的一部分\cite{rockoff2009}。
\end{pitfall}

桑顿对真实票据还有一条更根本的批评：一张真实票据背后有货物，一张“虚构的”票据背后也可能有一个资本雄厚、众所周知的出票人；认为“真实票据代表财产、虚构票据不代表，这种看法对一种票据多给了公道，对另一种少给了公道”\cite{thornton1802}。不过他也承认，虚构的票据一般更难按时付清，数量也更少受到限制：一个人实际卖出多少货物，多少给他的真实票据设了一个上限\cite{thornton1802}。这些批评与他关于利率的论证，留给\chref{ch:bullion}。

回到\secref{sec:merc-law}留下的问题：劳说他的纸币数量始终与需求相等，可需求是在什么价格、什么利率下的需求？真实票据论只回答了一半。它把一家银行的放款拴在真实的交易上，物价则由兑换金银来锚定；去掉兑换这个锚，按 Mints 的批评，票据的名义价值跟着物价上涨，“只贴现真实票据”就管不住货币的总量。

\section{艾尔银行}\label{sec:smith-ayr}

这一节回答：艾尔银行是怎样成长、怎样倒下的？斯密从中得出了什么？

\subsubsection{一家要解救国家的银行}
斯密说：“在这一片抱怨与困境之中，苏格兰设立了一家新银行，明说是为了解救国家的困境。”他的评价是：“用意是慷慨的，执行却是不审慎的”\cite{smith1776}。这家银行“似乎几乎不区分真实票据与循环票据，一律贴现”；它“公开的原则”是，只要有合理的担保，就预付改良土地之类回收最慢、最远的项目所需的全部资本\cite{smith1776}。按 Rockoff 文后的大事年表，道格拉斯—赫伦公司（艾尔银行）1769 年 11 月 6 日开业；Rockoff 说，到 1772 年，它发出的银行券占苏格兰各银行的 25\%，存款占 25\%，资产占 40\%，而两家公共银行合起来只占资产的 21\%；“三年之内，艾尔银行成了一个巨人”\cite{rockoff2009}。

它的资本并不充实。斯密说，两次认股共认缴资本十六万英镑，实缴只有八成\cite{smith1776}；“股东中很大一部分人在缴第一期股款时，就在本行开了现金账户”，董事们觉得应当像对待别人一样慷慨地对待自己的股东，就“允许他们中的许多人从这个现金账户借钱，来缴以后各期的股款”；斯密说，这样的缴款“只是把刚从一个金库里拿出来的钱放进另一个金库”，它的金库“从来没有充实过”\cite{smith1776}。银行的后盾是股东无限的责任：股东们的地产“值几百万”\cite{smith1776}。

\subsubsection{一本亏损的账}
斯密给出了这家银行停业时的几个数字：流通中的银行券“大约二十万英镑”；为了应付兑现，它不断在伦敦开汇票筹钱，停业时这些汇票的总额“超过六十万英镑”；它贷给各种人的钱“超过八十万英镑，利率 5\%”\cite{smith1776}。在伦敦开票筹钱，按\secref{sec:smith-realbills}的算法，连利息带佣金要付 8\% 以上；所以，斯密说，这家银行“在它全部业务的四分之三以上，损失超过 3\%”\cite{smith1776}。

\begin{example}{艾尔银行的损益}\label{ex:smith-ayr}
用斯密的数字，贷款 80 万英镑，利率 5\%；资金来源是 20 万英镑银行券与 60 万英镑在伦敦开出的汇票，后者的成本按 8\% 计。（1）一年的利息收入是多少？伦敦汇票的成本是多少？（2）靠伦敦汇票筹来的部分占贷款的几成？这部分一年亏多少？银行券那部分（暂不计准备与兑现的费用）赚多少？合计多少？（3）实缴资本约为多少？按（2）的亏损，几年就会亏掉？
\solution （1）$800{,}000\times5\%=40{,}000$ 英镑；$600{,}000\times8\%=48{,}000$ 英镑。（2）$600/800=75\%$，即斯密的“四分之三以上”；亏 $600{,}000\times(8\%-5\%)=18{,}000$ 英镑；银行券的部分赚 $200{,}000\times5\%=10{,}000$ 英镑；合计每年亏约 8,000 英镑，还没算收集金银的费用与坏账。（3）$160{,}000\times80\%=128{,}000$ 英镑，按每年 8,000 英镑，要 16 年才亏掉；可是实缴资本有很大一部分是向本行借来缴的，真正的缓冲要小得多，而贷给改良土地的钱短期内收不回来。银行不是被这笔利差拖垮的，是被不能兑现拖垮的。
\end{example}

\subsubsection{1772 年 6 月}
本章开头讲了危机的经过：6 月 8 日，那位合伙人的银行倒闭，按 Rockoff 的大事年表，它与艾尔银行“有大量的往来”；6 月 9 日他逃往法国；6 月 12 日消息传到爱丁堡，艾尔银行遭到挤兑，6 月 22 日停止兑付\cite{rockoff2009}。休谟在信里说的另一件事是苏格兰的卡伦公司（Carron Company）“摇摇欲坠，这是整场灾难里最大的灾难之一，因为它雇用了将近一万人”；苏格兰银行为它贴现票据，帮它渡过了危机\cite{rockoff2009}。艾尔银行向英格兰银行求援，Rockoff 称它为“那个时代的雷曼兄弟”；英格兰银行答应贷给它 30 万英镑，条件却苛刻到交易始终没有做成\cite{rockoff2009}。

危机随即传到欧洲大陆。Rockoff 说，1773 年 1 月的第一个星期，“伦敦与阿姆斯特丹之间的贸易与金融陷于停顿”；汉堡、斯德哥尔摩与圣彼得堡都受到了影响\cite{rockoff2009}。1773 年 1 月 10 日是个星期天，英格兰银行准许凭银行券与政府债券提取金银，据说一位荷兰银行家拿到了 50 万英镑（Rockoff 引 Kindleberger）\cite{rockoff2009}。同一个月，阿姆斯特丹市开办了一个由阿姆斯特丹银行（\secref{sec:finrev-amsterdam}）做后盾的贷款所；在圣彼得堡，叶卡捷琳娜二世（Catherine II，1762—1796 年在位）为英国商人提供了担保\cite{rockoff2009}。艾尔银行在 1772 年 9 月 28 日就已重新开门，承诺全额兑付银行券；1773 年 8 月 12 日，股东们决定清算\cite{rockoff2009}。

\subsubsection{斯密的判断}
斯密说，这家银行让那些投机者“比本来能做的多做了大约两年”；“如果他们中的大部分人早两年被迫停下来，对他们自己、对他们的债主、对他们的国家都会好得多”\cite{smith1776}。做循环汇票的人都转到这家新银行，在那里“受到张开双臂的欢迎”，别的银行于是“很容易地走出了那个致命的圈子”；所以，斯密说，从长远看，这家银行“加重了它想解救的国家的真实困境”，却把它想取代的对手“从极大的困境中解救了出来”\cite{smith1776}。他把这家银行比作一个水池，流出去的水没有流回来，只能“雇许多人不断提着水桶，到几英里外的井里打水来补充”；这个办法即使行得通、对银行有利，国家也得不到好处：它一点也不能增加可以出借的钱，“只会把这家银行变成整个国家的一种普通的贷款所”，而一家银行把钱借给也许五百个人，董事对其中大多数所知甚少，挑选借款人未必比一个私人更明智\cite{smith1776}。

\section{怎样管银行}\label{sec:smith-regulation}

这一节回答：斯密主张对银行业怎样管？这些规定在 1772 年起了作用吗？今天的人怎样看？

\subsubsection{小额银行券与防火墙}
斯密最具体的建议是禁止小额银行券。他说，“在伦敦，没有面额低于十英镑的银行券流通，纸币大体限于商人之间的流通”；在苏格兰，银行券小到“二十先令”，纸币就进入了商人与消费者之间的流通\cite{smith1776}。约克郡的一些纸币甚至小到六便士。小额银行券让小本的银行家也能发钞，“这样的穷银行家经常破产，可能给许多收了他们银行券的穷人带来很大的不便，有时甚至是很大的灾难”；所以，“也许最好是全国任何地方都不发行面额低于五英镑的银行券”\cite{smith1776}。

斯密承认，这样的限制，“是对天赋自由的明显侵犯，而法律的本分是保障这种自由，不是侵犯它”；可是，“为了防止火灾蔓延而规定修筑隔墙，同样是对天赋自由的侵犯，与这里提议的银行业管制完全属于同一类”\cite{smith1776}。他的条件很明确：“如果禁止银行家发行面额低于一定数额的流通银行券或见票即付的票据，如果他们有义务一经提示就立即无条件地兑付这些银行券，那么，在其他一切方面，他们的营业都可以完全自由，而不危及公众的安全。”\cite{smith1776}在这两条以内，银行越多越好：“最近英格兰与苏格兰两地银行公司的增多，许多人为之惊慌，其实不但没有减少，反而增加了公众的安全”；“把整个流通分成更多的部分，任何一家公司的倒闭，这种事在事物的进程中总会发生，对公众的影响就小了”；“这种自由竞争也迫使所有银行家对客户更宽厚，免得对手把客户抢走。”\cite{smith1776}

斯密还主张限制利率。在大不列颠，政府借钱付 3\%，有可靠担保的私人付 4\% 到 4.5\%，“现行的法定利率 5\% 也许和任何利率一样合适”；如果法定利率定到 8\% 或 10\%，“大部分要借出去的钱就会借给浪费者与投机者，只有他们愿意付这样高的利息”\cite{smith1776}。

\subsubsection{这些规定管用吗}
Rockoff 把斯密的建议归纳为四条：“（1）一部合理的高利贷法，（2）禁止银行券上的选择条款，（3）把真实票据原则作为银行业文化的一部分，（4）禁止面额低于五英镑的银行券”；他说，“前两条在 1772 年危机爆发时已经实行了”\cite{rockoff2009}。选择条款是斯密也讲到的一种做法：银行券上写明见票即付，或者“由董事选择，在提示后六个月”付款，外加这六个月的法定利息\cite{smith1776}。按 Rockoff 的大事年表，1766 年 5 月 15 日起禁止苏格兰银行券使用选择条款，6 月 1 日起禁止面额低于 1 英镑的银行券\cite{rockoff2009}。“尽管斯密主张，五英镑的下限无论在法律上还是在实践上，苏格兰都从未采用”；艾尔银行发行过 1 英镑与 1 基尼（1.05 英镑）的银行券\cite{rockoff2009}。Rockoff 认为，如果五英镑的下限在 1769 年以前就实行，“可能会有相当大的作用”，因为艾尔银行就得“靠范围更窄、更懂行的一类持券人”筹一部分钱\cite{rockoff2009}。

可是，Rockoff 的结论是：“艾尔银行的倒闭不是因为不懂审慎经营的规则，而是因为决意不理会它们”；“主要的问题不是复杂的失误，而是过度的内部人贷款”\cite{rockoff2009}。股东借本行的钱来缴股款，正是斯密自己记下的。

\subsubsection{国家的大机器}
斯密对英格兰银行另眼相看：它“不仅是一家普通的银行，而且是国家的一部大机器”；它为商人贴现票据，好几次“支撑了不仅英格兰、而且汉堡与荷兰的主要商号的信用”；1763 年有一次，它据说为此“在一个星期里”预付了约 160 万英镑，其中很大一部分是金块，“不过，我不敢担保数额有这么大，时间有这么短”\cite{smith1776}。Rockoff 说，英格兰银行在 1763 年的危机里充当了最后贷款人，但那场危机的中心在阿姆斯特丹，对斯密关于银行的想法似乎没有大的影响\cite{rockoff2009}。Rockoff 说，英国银行危机的真正结束，是英格兰银行承担起\textbf{最后贷款人}的角色；他写于 2009 年，在这句话里加了“直到现在”，指的是刚刚发生的 2008 年危机。“斯密没有明确赞成用这些权力来缓和或防止危机”\cite{rockoff2009}。他引 Schwartz（1987）的话说，英格兰银行采取了最后贷款人的政策以后，1866 年以后英格兰再没有发生银行危机\cite{rockoff2009}。最后贷款人的原则通常与沃尔特·白芝浩（Walter Bagehot，1826—1877）的名字连在一起，Rockoff 说白芝浩的看法比通常的描述复杂，见\chref{ch:goldstandard}；Rockoff 说，白芝浩自己的偏好，如果能从头来过，是“苏格兰那样的自由银行制度”\cite{rockoff2009}。

\begin{compare}{隔墙与消防队}{斯密的隔墙}{最后贷款人}
\cmprow{思路}{事先限制银行，让一家倒闭不连累别家}{危机发生时由中央银行放款，制止挤兑}
\cmprow{主要手段}{禁止小额银行券；见票即付；法定利率；银行越多越好}{对有偿付能力的银行放款（\chref{ch:goldstandard}）}
\cmprow{1772 年的情形}{前两条已经实行，五英镑的下限没有；艾尔银行不理会规则}{英格兰银行的条件苛刻，没有贷给艾尔银行}
\cmprow{英格兰的结果}{——}{按 Schwartz（1987）的说法，1866 年以后没有银行危机；2008 年又有了（\chref{ch:gfc}）}
\cmprow{代价}{限制了小额信用与“天赋自由”}{银行敢于冒险，因为知道有人兜底}
\cmpsum{斯密想靠规则与分散来防火，后来的英国靠中央银行来救火；两者并不互相排斥，今天的银行监管与最后贷款人是同时存在的。}
\end{compare}

\begin{pitfall}[斯密主张银行业完全放任]
斯密在多数行业信任市场，对银行却主张限制面额、强制兑付与限制利率；Rockoff 说，有些研究斯密的人，比如 West，为这种表面上的矛盾而困惑\cite{rockoff2009}。Rockoff 的解释是：“斯密是一个经验主义者，愿意根据当时与历史的经验修改自己的看法。”\cite{rockoff2009}斯密的“完全自由”是有条件的，条件就是那两条规定；他用隔墙作比，正是说某些对自由的限制本身是为了保护自由。
\end{pitfall}

\section{进步、停滞与衰退}\label{sec:smith-states}

这一节回答：斯密怎样区分进步、停滞与衰退的国家？他的证据是什么？哪些判断经得起检验？

\subsubsection{工资看的是增长}
斯密说：“引起劳动工资上涨的，不是国民财富的实际大小，而是它的不断增长。”\cite{smith1776}所以，“英格兰今天无疑比北美的任何地方都富得多”，北美的劳动工资“却比英格兰的任何地方都高得多”\cite{smith1776}。他引纽约省的工资：1773 年，“在最近的骚乱开始以前”，普通劳工每天挣“三先令六便士殖民地货币，合两先令英镑”；造船木匠“十先令六便士殖民地货币，外加一品脱朗姆酒，值六便士英镑，合计六先令六便士英镑”；房屋木匠与砖瓦匠“八先令殖民地货币，合四先令六便士英镑”；缝纫帮工“五先令殖民地货币，合大约两先令十便士英镑”；“这些价格都高于伦敦的价格”\cite{smith1776}。

\begin{example}{两种先令}\label{ex:smith-wages}
（1）由普通劳工的工资，纽约的殖民地货币一先令值多少英镑的先令？一英镑的先令折多少殖民地先令？（2）用这个比价验算造船木匠、房屋木匠与缝纫帮工的英镑工资。（3）哪一个数与比价对不上？差多少？
\solution （1）$3.5/2=1.75$ 殖民地先令折 1 英镑先令，殖民地先令值 $1/1.75\approx0.571$ 英镑先令。（2）造船木匠 $10.5/1.75=6$ 先令，加朗姆酒 6 便士，合 6 先令 6 便士，对得上；缝纫帮工 $5/1.75\approx2.857$ 先令，即 2 先令 10 便士多一点，斯密写作“大约”；房屋木匠 $8/1.75\approx4.571$ 先令，即 4 先令 7 便士不到。（3）房屋木匠：斯密写作 4 先令 6 便士，按比价应为约 4 先令 6.9 便士，差不到 1 便士；他大概是取了整数。这里的比价就是\chref{ch:external}所说的汇率：同一个“先令”，在两个地方是两种钱。
\end{example}

\subsubsection{北美：靠的是繁殖}
斯密说，北美的人口二十或二十五年翻一番，而且“这种增长主要不是由于不断输入新居民，而是由于人类自身的大量繁殖”\cite{smith1776}。劳动的需求决定人口的增减：斯密认为，这种需求在北美、欧洲与中国各不相同，“使人口的增殖在第一个地方迅速前进，在第二个地方缓慢渐进，在最后一个地方完全停滞”；他说，劳动贫民的处境，“在停滞的状态下是艰难的，在衰退的状态下是悲惨的”\cite{smith1776}。

\subsubsection{中国：法律与制度}
斯密对中国的判断有三层。第一，工资低：许多旅行者的记述“都说中国的劳动工资低，劳动者难以养家”；“在广州附近，有好几百户人家，一般的说法是好几千户，在陆地上没有住处，常年住在江河运河的小渔船上”\cite{smith1776}。第二，中国“也许停滞不前，却似乎没有退步”\cite{smith1776}。第三，停滞的原因是制度：“中国似乎长期停滞，大概早就积累到了与它的法律与制度的性质相适应的全部财富”；“一个忽视或者鄙视对外贸易、只许外国船只进入一两个港口的国家，不能做它在不同的法律与制度下所能做的那么多的生意”；在那里，富人与大资本的所有者有相当的安全，穷人与小资本的所有者却“几乎没有任何安全”，随时可能“在司法的名义下被下级官吏掠夺”；所以，“据说中国普通的利率是 12\%”\cite{smith1776}。

斯密没有中国的人口数字。《清史稿》记载，户部登记的民数，乾隆二十九年（1764）为“二万五百五十九万一千一十七口”，乾隆六十年（1795）为“二万九千六百九十六万五百四十五口”，并说“经累朝休养生息，故户口之数，岁有加增”\cite{qingshigao}。这是\secref{sec:smith-test}的一个反例。

\subsubsection{孟加拉：一家压迫人的公司}
斯密说，英国宪政保护与统治的北美，与“在东印度压迫人、专横跋扈的那家商业公司”治下的地方，两者精神的差别，“也许没有比这两个地方的不同状况更好的说明了”\cite{smith1776}。在孟加拉，“钱常常以 40\%、50\% 与 60\% 的利率借给农民，下一季的收成抵押来还债”\cite{smith1776}。关于饥荒，斯密在第四篇第五章说：“几年前孟加拉的干旱，也许本来会造成一场很严重的歉收。东印度公司的雇员对大米贸易施加的一些不当的规定、一些不明智的限制，也许促成了把歉收变成饥荒。”\cite{smith1776}他用了两个“也许”。

\begin{case}{孟加拉 · 1770—1772 · 饥荒与田赋}
亨特（W. W. Hunter，1840—1900）1868 年出版的《孟加拉农村编年史》，第二章写 1770 年的饥荒。他说：“1770 年 5 月底以前，官方估计人口的三分之一已经消失”；这个估计“为所有官方的与最严谨的”非官方作者所接受\cite{hunter1868}。他自己由此推算，饥荒“在九个月里夺去了一千万人的生命”\cite{hunter1868}。这是亨特从官方的三分之一与他对当时人口的估计推出来的，不是当时的统计；斯密的数字是“一年里饿死三四十万人”，两者相差二十五到三十三倍；斯密的这句话在 1776 年的初版里就有\cite{hunter1868,smith1776}。

人口少了三分之一（死去的与逃走的），田赋却没有少。亨特在书后附了一封信：“1772 年 11 月 3 日，于威廉堡”，沃伦·黑斯廷斯（Warren Hastings，1732—1818）与他的参事会写给东印度公司董事会\cite{hunter1868}。信里说，“尽管全省至少失去了三分之一的居民，耕地也随之减少”，1771 年的净收入“甚至超过了 1768 年”\cite{hunter1868}。信后的账目摘要显示，1768—1769 年的净收入是 1,52,54,856 西卡卢比，1771—1772 年扣除缺额以后是 1,53,33,660 卢比（印度的写法，即 15,254,856 与 15,333,660）\cite{hunter1868}。信里自己解释了原因：收入没有下降，“是因为它被强行维持在原来的标准上”；其中一项叫“纳贾伊”（Najay），“是向各类田地上实际还活着的居民加征的一种摊派，用来弥补邻居死去或逃走而损失的地租”\cite{hunter1868}。

饥荒为什么这样严重，当时就有两种说法。一种说法怪罪对粮食贸易的干预。斯密说，即使在种稻的国家，干旱也很少普遍到必然造成饥荒的程度，“只要政府允许自由贸易”；政府为了缓解歉收，命令所有的商人按它认为合理的价格卖粮，就会阻止他们把粮食运到市场上，“有时在一季开始的时候就造成饥荒”\cite{smith1776}。亨特同意这种看法。他是拿 1866 年奥里萨饥荒的经验来读 1770 年的：他推断两次的天灾一样重，因为在条件相似的地方，“米价所反映的实际压力，两次饥荒完全一样”\cite{hunter1868}。加尔各答的参事会很早就“发布公告，禁止囤积或垄断粮食”；他说，“1770 年，政府禁止它所谓的粮食垄断，阻止了价格立即涨到它自然的水平”；如果让私人去储存，粮价一开始就会上涨，人们从一开始就会少吃，存粮就能“省着用，把压力平均分摊到整整九个月”；结果，粮价在饥荒的头几个月一直是每磅四分之三便士，没有像 1865—1866 年那样立即涨到一个半便士\cite{hunter1868}。另一种说法怪罪公司的雇员自己。亨特书后的附录收了董事会事后的一封信：信里引用的指控说，英国绅士们的代理人“不仅垄断粮食，而且强迫穷苦的农民卖掉下一季必需的种子”；董事会说，农民被迫把米卖给“这些垄断的欧洲人”，“我们有理由怀疑，他们只能是我们公司里有些地位的人”\cite{hunter1868}。亨特也记下，董事会对参事会的救济措施似乎很满意，却责备它“纵容了粮食的垄断与私人贸易”\cite{hunter1868}。两种说法并不互相排斥：禁止商人囤粮与公司雇员自己垄断粮食，可以同时发生。

斯密说的是：孟加拉在衰退，维持劳动者的基金在缩小；造成这种状况的是一家压迫人的公司；公司雇员对米粮贸易的不当规定与限制“也许”把歉收变成了饥荒。他没有说饥荒是公司造成的，也没有引用孟加拉的田赋数字。黑斯廷斯的信说明了另一件事：死者的税由生者来交。
\end{case}

表~\ref{tab:smith-states} 把斯密的三种国家与本章的数据放在一起。

\begin{table}[htbp]\centering\small
\caption{斯密的进步、停滞与衰退}\label{tab:smith-states}
\begin{tabularx}{\linewidth}{@{}>{\raggedright\arraybackslash}p{0.12\linewidth}YYYY@{}}
\toprule
 & 北美殖民地 & 英国（欧洲） & 中国 & 孟加拉 \\\midrule
斯密的判断 & 迅速进步 & 缓慢渐进 & 停滞，不退步 & 衰退 \\
斯密的证据 & 工资高于伦敦；20 或 25 年翻一番 & 不到 500 年不翻一番 & 工资低，广州的船户；利率 12\% & 一年饿死三四十万人；利率 40\%—60\% \\
斯密说的原因 & 英国的宪政 & —— & 法律与制度，只开一两个港口 & 压迫人的商业公司 \\
本章的数据 & 1700—1770 年，22.6 年翻一番 & 英格兰 1700—1770 年 205 年、1750—1780 年 107 年翻一番 & 登记民数 1764—1795 年增加 44.4\%，58 年翻一番 & 1770 年官方估计人口消失三分之一；1771 年的净收入超过 1768 年 \\
\bottomrule
\end{tabularx}
\datasource{斯密《国富论》第一篇第八、九章与第四篇第五章；北美据美国人口普查局（1975）系列 Z 1；英格兰据英格兰银行《千年宏观数据》表 A2（Broadberry 等 2015）；中国据《清史稿》卷一百二十；孟加拉据亨特（1868）及其附录 A 所载 1772 年 11 月 3 日的信。翻一番的年数按两年之间的平均增长率算，由本书计算。}
\end{table}

\section{数据检验与反例}\label{sec:smith-test}

这一节回答：斯密的人口翻番之说对不对？公家的浪费是否让英国停了下来？苏格兰的纸币是否像他说的那样过多？艾尔银行为什么倒？

\subsubsection{翻一番的年数}
图~\ref{fig:smith_population} 画出了北美殖民地与英格兰的人口。北美殖民地 1700—1770 年增加到 8.56 倍，平均 22.6 年翻一番，正在斯密说的二十到二十五年之间；1700—1780 年是 23.1 年\cite{census1975}。普查局的数字包括黑人：1770 年有 459,822 人，占 21.4\%。斯密自己说，在以谷物为主的英国殖民地，“绝大部分活是自由人做的”；在甘蔗殖民地，“全部的活是奴隶做的”，在烟草殖民地也有很大一部分\cite{smith1776}。这些奴隶本人或先辈是从非洲被强行贩运来的，干活没有工资；斯密说北美的劳动工资高，说的是自由劳动者。按序列说明，17 世纪黑人的数字“无疑偏低”，是根据买卖的记载、有关奴隶制的法律、偶尔也根据大体准确的报告估计出来的\cite{census1975}。只算白人，1700—1770 年也是 24.0 年翻一番，仍在斯密的范围之内；黑人人口翻一番只要 17.3 年。翻一番的年数只检验增长的速度，分不出其中有多少来自出生、多少来自移民与奴隶贸易；斯密说的“不是由于输入新居民”，对黑人尤其未必成立。

斯密转述的通行说法，欧洲的人口“不到五百年不会翻一番”，对英格兰太悲观了。按千年数据，英格兰的人口 1700 年约 520 万，1770 年约 658 万，平均约 205 年翻一番；到 1750—1780 年，已经加快到约 107 年\cite{boe2017,broadberry2015}。斯密写书的时候，英格兰正在从“缓慢渐进”的状态里走出来，这是\chref{ch:divergence}的题目。

\fig{smith_population}{北美、英格兰与清代的人口增长。(a) 看各条线的斜率：纵轴为对数刻度，直线越陡，增长越快。实线是北美殖民地的人口（白人与黑人，每十年一个数，1630—1780），褐色虚线是其中的黑人，浅色线是英格兰的人口（1660—1780，逐年）。(b) 看柱子与斯密的说法差多少：翻一番所需的年数，按两年之间的平均增长率计算；浅色带是斯密说的北美（20—25 年），竖线是他说的欧洲（至少 500 年）。数据：美国人口普查局（1975）系列 Z 1；英格兰银行《千年宏观数据》表 A2；《清史稿》卷一百二十。}{fig:smith_population}

清代是一个反例。斯密说中国停滞，又说劳动的需求决定人口的增殖；按他自己的标准，停滞的国家人口不应增长。可是《清史稿》的登记民数，1764—1795 年增加了 44.4\%，平均每年 1.19\%，58 年翻一番，比同期的英格兰还快\cite{qingshigao}。这个反例有它的限度：这是户部的登记数，不是普查，登记的办法也可能变；人口增长也不等于工资上升，斯密说的低工资与人口增长可以同时存在。但斯密“停滞”的判断，至少要限于人均的意义上。

\subsubsection{公家的浪费与人均产出}
斯密说政府的浪费延缓了进步，却没能让它停下来。图~\ref{fig:smith_debt} 支持后一半：1700—1776 年，国债增加到约 9.7 倍，人均实际 GDP 仍然增加了 22.8\%\cite{boe2017}。前一半，没有战争英国会长得多快，数据回答不了，因为没有一个不打仗的英国可以比较。

\subsubsection{苏格兰的纸币多不多}
Rockoff 的表 1 据 Checkland 给出了苏格兰银行的资产负债表\cite{rockoff2009}。1753 年，它的银行券是 6.4 万英镑，存款 3.1 万英镑；1764 年，银行券 16.3 万英镑，存款 10.2 万英镑。Rockoff 说：“银行券增加到 2.5 倍以上，年率 8.5\%；存款增加到将近 3.3 倍，年率 10.8\%。”\cite{rockoff2009}这两个年率是对数增长率，$\ln(163/64)/11=8.5\%$，按复利算是 8.9\%；存款的复利年率是 11.4\%。斯密说苏格兰的纸币二十五年多以来一直“略多于”流通能吸收的数量，Rockoff 也说，“小额银行券狂热”在 1763 年与 1764 年达到顶峰\cite{rockoff2009}。但他对原因有另一种猜测：“我的猜测是，小面额铸币的短缺是最重要的因素”；“七年战争（1754—1763）造成了英国铸币的净流失”\cite{rockoff2009}。按这种看法，小额银行券是在补铸币的缺，不全是投机者的狂热。

\subsubsection{艾尔银行为什么倒}
斯密把艾尔银行的倒闭归于不区分真实票据与循环票据、把短期的钱贷给回收很慢的土地改良。Rockoff 同意它违反了斯密的规则，但认为原因不在于不懂规则：股东借本行的钱缴股款，银行把钱贷给自己人，是“过度的内部人贷款”\cite{rockoff2009}。两种解释都能在斯密自己的记述里找到根据。

\begin{deeper}[版本与史料的出入]
\textbf{初版与后来的版本。}古腾堡的文本出自第三版或更晚的某一版。本章引用的主要段落在 1776 年初版的扫描本里都能找到：别针工场的“四万八千多枚”、孟加拉的“三四十万人”、纽约的工资、孟加拉 60\% 的利率、纸币节省的“五分之一”、“真实的债权人……真实的债务人”、空中的车道与“代达罗斯的翅膀”、隔墙、“完全自由”与“超过 1.45 亿英镑”（初版用文字写作“一亿四千五百多万”）\cite{smith1776}。有两处不同：\exref{ex:smith-redraw} 第（2）问引的那条注，说靠流通筹钱的另一种做法“至少要花 14\%”，初版没有，初版第 375—377 页只有正文的 8\%；第四篇第八章的名句“消费是一切生产的唯一目的”，所在的那一章（“重商主义的结论”）初版没有，初版第四篇的第八章是后来的第九章，讲重农主义\cite{smith1776}；本章不用这句话。古腾堡本把“养活许多仆人”的 of 印成了 or，本章照原文的意思翻译。

\textbf{1772 年的日期。}Rockoff 的正文说艾尔银行 6 月 22 日“被迫停止兑付银行券”，他文后的大事年表说 6 月 25 日“停止以金银兑付银行券”；本章用正文的日期，并在这里注明。逃往法国的那位合伙人，Rockoff 说是伦敦的投机者亚历山大·福代斯（Alexander Fordyce），“关键人物之一”\cite{rockoff2009}；我们没有找到可以核对的生卒年，所以正文里不写他的名字。

\textbf{英格兰银行的 160 万英镑。}斯密自己说他“不敢担保”数额与时间。

\textbf{孟加拉的死亡。}亨特的“一千万人”是他根据官方的三分之一推算的，斯密的“三四十万人”没有说明来源；两者都不是统计。黑斯廷斯的信说“至少三分之一的居民”，亨特引的官方估计是 1770 年 5 月底以前的“三分之一”。亨特书的识别文本把信里的年份“1771”读成了“177 1”。

\textbf{殖民地的人口。}普查局的序列由 Sutherland 汇编，原始数据来自殖民地官员向贸易与种植园委员会的报告，只分白人与黑人两类，原住民基本不在其中，只有 1780 年罗得岛与康涅狄格两个格子的注说“包括一些印第安人”；本章按页面图像读出\cite{census1975}。序列说明还说，一些没有交代依据的估计，是按人口增长的总体趋势挑选进来的，所以用它检验增长速度并不完全独立。

\textbf{《清史稿》的数字。}“二万五百五十九万一千一十七口”里的“万”用了两次，读作 $(2\times10{,}000+559)\times10{,}000+1{,}017=205{,}591{,}017$。

\textbf{Rockoff 的年率。}Rockoff 的 8.5\% 与 10.8\% 是对数增长率，不是复利年率（见上文）。
\end{deeper}

\section{今天还剩什么}\label{sec:smith-legacy}

斯密的许多判断后来被修正了：纸币不会抬高物价的说法，桑顿指出它忽略了流通的快慢，在不兑现的纸币下更站不住（\chref{ch:bullion}）；中国停滞的判断，至少在人口上不对；“节俭必然有益”也受到了质疑。另外几个问题今天还在。

\textbf{增长从哪里来。}式~\eqref{eq:smith-twofactors} 把人均产出分成生产率与从业比例；斯密把生产率的提高归于分工、机器与资本积累。把增长分解成资本、劳动与技术的贡献，是\chref{ch:solow}的增长核算；英国为什么在 18 世纪末加速，是\chref{ch:divergence}的题目。

\textbf{储蓄与需求。}“省下的同样被消费掉”，经过萨伊定律（\chref{ch:say}），成了凯恩斯要推翻的对象：一个人节俭是美德，大家同时节俭，需求就可能不足（\chref{ch:generaltheory}）。

\textbf{银行发钱的锚。}真实票据论在金银本位下是一条审慎经营的规则，在不兑现的纸币下就失去了锚；桑顿的批评与英格兰银行停止兑付以后的争论，见\chref{ch:bullion}。

\textbf{防火与救火。}斯密的隔墙与白芝浩的最后贷款人，今天都还在：资本与流动性的规定是隔墙，中央银行的紧急贷款是救火（\chref{ch:goldstandard}）。艾尔银行的内部人贷款与“那个时代的雷曼兄弟”，在\chref{ch:gfc}的全球金融危机里有它的回声。

\textbf{钱只是一个轮子。}金银是“死的资财”，节省下来的金银能变成资本。斯密的银行留五分之一的金银、发出五倍于准备的银行券，与\secref{sec:money-creation}里准备与货币之比的算术相同；不同的是约束：斯密的银行券随时要兑换金银，多发的部分会回来兑现（式~\eqref{eq:smith-reflux}），今天的银行发放贷款就创造存款，存款兑换的是中央银行的准备金，约束它的是利率、竞争与监管（\secref{sec:money-creation}）。

\begin{takeaways}
\begin{itemize}
\item 斯密说，国民财富是一年的劳动产出，不是金银。人均产出等于劳动生产率乘以从业比例，$Y/N=(Y/L)(L/N)$，式~\eqref{eq:smith-twofactors}；比例有上限，生产率没有。别针工场里十个人一天做四万八千多枚，每人 4,800 枚，是一个人单干的 240 倍以上。分工受市场范围的限制，市场范围取决于运输：一条船运 200 吨，抵得上 50 辆大车、100 个人与 400 匹马。
\item 资本靠节俭积累，“节俭，而不是勤劳，是资本增加的直接原因”；生产性劳动是固定在可以出售的物品上的劳动。斯密认为省下的钱照样被消费，只是换了消费者，这是后来“普遍过剩”之争的起点。英国的国债从 1700 年的 1,410 万英镑增加到 1776 年的 1.366 亿英镑，人均产出仍然增加了 22.8\%。
\item 纸币代替金银，省下的金银可以换成资本：$\Delta Y/Y=(M/Y)(1-\rho)r$，式~\eqref{eq:smith-saving}；按 Rockoff 的计算，只有产出的 0.27\% 到 1.60\%。纸币随时可兑现时，多发的部分会回来兑现，银行要按一比一准备，$R=\rho\bar N+(N-\bar N)$，式~\eqref{eq:smith-reflux}。
\item 真实票据论：银行只应贴现真实的、到期会付清的商业票据，不应贷出固定资本。循环开票的成本 $c=i+nk$，式~\eqref{eq:smith-redraw}，超过 8\%，而利润只有 6\% 到 10\%。在金银本位下这是一条审慎经营的规则；在不兑现的纸币下，票据的名义价值随物价上涨，它就控制不了货币的数量（Mints）。
\item 艾尔银行 1769 年开业，三年里占到苏格兰各银行资产的 40\%，1772 年 6 月停止兑付：贷出 80 万英镑，利率 5\%，其中 60 万靠成本超过 8\% 的伦敦汇票筹来。斯密主张禁止五英镑以下的银行券，条件以内银行业“完全自由”；Rockoff 认为艾尔银行倒于内部人贷款，后来平息英国银行危机的是英格兰银行做最后贷款人，直到 2008 年的危机。
\item 斯密把北美、中国与孟加拉看作进步、停滞与衰退的三种国家。北美殖民地的人口 22.6 年翻一番，与他说的相符；英格兰的人口增长比他说的快，清代的登记民数 1764—1795 年增加了 44.4\%；孟加拉 1770 年的饥荒之后，田赋被“强行维持在原来的标准上”。
\end{itemize}
\end{takeaways}

\begin{selfcheck}
\begin{enumerate}
\item 一个国家的人口为 1,000 万，从事有用劳动的有 450 万，一年的产出为 9 亿英镑。按式~\eqref{eq:smith-twofactors}，劳动生产率与人均产出各是多少？如果从业的人增加到 500 万，生产率不变，人均产出增加多少？
\item 一条船由 8 个人驾驶，运 200 吨货物。改用斯密的大车，运同样的货物要多少辆车、多少个人？如果运河让一条 6 个人的船也能走到内地，用人减少多少？
\item 一个国家流通的金银为 600 万英镑，一年的产出为 6,000 万英镑。银行用纸币代替金银，留 1/4 的准备，换回来的资本利润率为 8\%。按式~\eqref{eq:smith-saving}，产出增加多少？
\item 一家银行的流通能吸收 6 万英镑的银行券，它留 1.2 万英镑金银。如果它发出 7 万英镑，按式~\eqref{eq:smith-reflux} 要留多少金银？比按同一比例算多留多少？
\item 一个商人每三个月开一次汇票续借，每次佣金 0.75\%，利率 5\%。按式~\eqref{eq:smith-redraw}，一年的成本是多少？如果他的项目利润率是 8\%，划算吗？
\item 一家银行贷出 100 万英镑，利率 5\%；资金中 30 万是银行券，70 万是成本 7\% 的借款。一年的损益是多少？借款占资金的几成？
\end{enumerate}
\end{selfcheck}

\begin{answers}
\item 劳动生产率 $9\text{亿}/450\text{万}=200$ 英镑；人均产出 $9\text{亿}/1{,}000\text{万}=90$ 英镑，等于 $200\times0.45$。从业的人增加到 500 万，人均产出 $200\times0.5=100$ 英镑，增加 $11.1\%$。
\item $200/4=50$ 辆，$50\times2=100$ 个人；改用 6 个人的船，比 8 个人的船少 2 个人，即 25\%，比陆运少 94 个人。
\item $\Delta Y/Y=(6/60)\times(3/4)\times8\%=0.6\%$，即 36 万英镑。
\item $\rho=12{,}000/60{,}000=0.2$，$R=0.2\times60{,}000+10{,}000=22{,}000$ 英镑；按同一比例只要 $0.2\times70{,}000=14{,}000$ 英镑，多留 8,000 英镑。
\item $c=5\%+4\times0.75\%=8\%$；利润率 8\% 刚好抵消成本，没有剩余，还要承担项目失败与续借不成的风险，不划算。
\item 收入 $1{,}000{,}000\times5\%=50{,}000$ 英镑；借款成本 $700{,}000\times7\%=49{,}000$ 英镑；不计准备与兑现的费用，赚 1,000 英镑。借款占七成，这一部分每年亏 $700{,}000\times2\%=14{,}000$ 英镑，靠银行券的 15,000 英镑补回来。
\end{answers}

斯密的纸币有一个前提：随时可以兑换金银。如果这个前提不在了，纸币多了会不会抬高物价？银行只贴现真实票据，还能不能保证纸币不会多发？下一章去看英格兰银行停止兑付以后的那场争论。
